% source: erdos-872/phase4/theorem5_disjoint_carriers_note.md
% source: erdos-872/phase4/theorem6_rank3_squarefree_note.md

\section{Restricted Carrier Classes}\label{sec:restricted-class}

This section records two positive theorems showing that several natural
Prolonger classes are not genuinely hard.  In both cases Shortener resolves all
small-prime structure after only $O_\alpha(n/\log n)$ moves, leaving at most
$O_\alpha(n/\log n)$ total legal play.

\subsection{Disjoint small-prime carriers}

\begin{theorem}[Theorem 5]\label{thm:theorem5}
Fix $\alpha$ with $1/3<\alpha<1/2$, and set $y=n^\alpha$.  Suppose every
composite Prolonger move has all prime factors at most $y$, and suppose further
that distinct composite Prolonger moves have pairwise-disjoint prime supports.
Then Shortener has an online strategy forcing
\[
  \L(n)=O_\alpha\!\left(\frac{n}{\log n}\right).
\]
\end{theorem}

\begin{proof}[Proof sketch]
Let $B$ be the set of small primes that occur in some composite Prolonger move.
For each $p\in B$, let $C(p)$ denote its unique carrier and let $e(p)=v_p(C(p))$.

Shortener uses a three-phase priority strategy:
\begin{enumerate}
  \item whenever a legal prime exists, play a legal prime;
  \item if $C(p)$ has at least two distinct prime factors, earmark the repair
    move $p^{e(p)+1}$ and play it whenever legal and higher-priority moves are
    unavailable;
  \item for distinct carrier primes $p,q\in B$ with $C(p)\neq C(q)$, earmark
    the semiprime $pq$ and play it whenever legal and higher-priority moves are
    unavailable.
\end{enumerate}

The disjoint-support hypothesis guarantees legality: a power repair
$p^{e(p)+1}$ is incomparable with its own carrier because that carrier contains
another prime, and incomparable with every other carrier because no other
carrier uses $p$; similarly, a cross-carrier pair $pq$ is incomparable with all
earlier phase-2 powers and all composite carriers.

Once these repairs are exhausted, every prime $p\le y$ is resolved.  Indeed, if
some legal move $x$ were still divisible by $p$, then every prime factor of $x$
would lie in $B$, and either $x$ would involve primes from two distinct
carriers, in which case a phase-3 pair would divide $x$, or all prime factors
of $x$ would lie inside one carrier, in which case either $x$ is comparable with
that carrier or some phase-2 power repair divides $x$.  Both cases are
impossible.

Once every $p\le y$ is resolved, every future legal move lies in the sparse set
\[
  \mathcal S_y := \{x\le n : P^-(x)>y,\ \Omega(x)\le 2\},
\]
whose size is $O_\alpha(n/\log n)$.  The total move count is therefore
$O_\alpha(n/\log n)$.
\end{proof}

\subsection{Squarefree rank-\texorpdfstring{$\le 3$}{<=3} carriers}

\begin{theorem}[Theorem 6]\label{thm:theorem6}
Fix $\alpha$ with $1/3<\alpha<1/2$, and set $y=n^\alpha$.  Suppose every
composite Prolonger move is squarefree, supported on primes at most $y$, and of
rank at most $3$.  Then Shortener has an online strategy forcing
\[
  \L(n)=O_\alpha\!\left(\frac{n}{\log n}\right).
\]
\end{theorem}

\begin{proof}[Proof sketch]
Let $B$ be the set of carrier primes.  Shortener again uses a priority strategy:
\begin{enumerate}
  \item legal primes;
  \item squares $p^2$ for primes $p\in B$ not already played;
  \item squarefree semiprimes $pq$ on $B$;
  \item squarefree triples $pqr$ on $B$.
\end{enumerate}

After phases~1 and~2, every legal survivor is squarefree and supported on $B$.
Indeed, any prime divisor outside $B$ would already have been resolved by a
phase-1 prime move, and any repeated prime factor $p^2$ would have been removed
by a phase-2 square.  From there the phase-3 and phase-4 priorities eliminate
all rank-$2$ and rank-$3$ squarefree survivors directly.

The only nontrivial case is a surviving squarefree integer $x$ with at least
four prime factors from $B$.  Choose three of them, say $p,q,r$.  Then
$pqr\mid x$.  If the triple $pqr$ was legal after phase~3, phase~4 would have
played it.  If it was not legal, then some earlier move comparable with $pqr$
already existed.  Because all carriers are squarefree of rank at most $3$, any
such earlier comparable move must itself divide $x$; in every case $x$ is
therefore illegal.  So no legal move survives the four phases.

Finally, the number of played phase-1, phase-2, and phase-3 moves is
$O_\alpha(n/\log n)$ by trivial prime and semiprime counts, while the phase-4
triple count is also $O_\alpha(n/\log n)$ for each fixed $\alpha<1/2$.  Hence
the total game length is $O_\alpha(n/\log n)$.
\end{proof}

\begin{remark}
These theorems explain why the later universal carrier-mass obstruction must be
read correctly.  The block-product family is an obstruction to a \emph{proof
method}, not to the actual existence of strong Shortener strategies.  In fact,
block-product play lies inside the scope of Theorem~\ref{thm:theorem5}.
\end{remark}
